% Options for packages loaded elsewhere
\PassOptionsToPackage{unicode}{hyperref}
\PassOptionsToPackage{hyphens}{url}
\documentclass[
]{article}
\usepackage{xcolor}
\usepackage[margin=1in]{geometry}
\usepackage{amsmath,amssymb}
\setcounter{secnumdepth}{-\maxdimen} % remove section numbering
\usepackage{iftex}
\ifPDFTeX
  \usepackage[T1]{fontenc}
  \usepackage[utf8]{inputenc}
  \usepackage{textcomp} % provide euro and other symbols
\else % if luatex or xetex
  \usepackage{unicode-math} % this also loads fontspec
  \defaultfontfeatures{Scale=MatchLowercase}
  \defaultfontfeatures[\rmfamily]{Ligatures=TeX,Scale=1}
\fi
\usepackage{lmodern}
\ifPDFTeX\else
  % xetex/luatex font selection
\fi
% Use upquote if available, for straight quotes in verbatim environments
\IfFileExists{upquote.sty}{\usepackage{upquote}}{}
\IfFileExists{microtype.sty}{% use microtype if available
  \usepackage[]{microtype}
  \UseMicrotypeSet[protrusion]{basicmath} % disable protrusion for tt fonts
}{}
\makeatletter
\@ifundefined{KOMAClassName}{% if non-KOMA class
  \IfFileExists{parskip.sty}{%
    \usepackage{parskip}
  }{% else
    \setlength{\parindent}{0pt}
    \setlength{\parskip}{6pt plus 2pt minus 1pt}}
}{% if KOMA class
  \KOMAoptions{parskip=half}}
\makeatother
\usepackage{longtable,booktabs,array}
\newcounter{none} % for unnumbered tables
\usepackage{calc} % for calculating minipage widths
% Correct order of tables after \paragraph or \subparagraph
\usepackage{etoolbox}
\makeatletter
\patchcmd\longtable{\par}{\if@noskipsec\mbox{}\fi\par}{}{}
\makeatother
% Allow footnotes in longtable head/foot
\IfFileExists{footnotehyper.sty}{\usepackage{footnotehyper}}{\usepackage{footnote}}
\makesavenoteenv{longtable}
\usepackage{graphicx}
\makeatletter
\newsavebox\pandoc@box
\newcommand*\pandocbounded[1]{% scales image to fit in text height/width
  \sbox\pandoc@box{#1}%
  \Gscale@div\@tempa{\textheight}{\dimexpr\ht\pandoc@box+\dp\pandoc@box\relax}%
  \Gscale@div\@tempb{\linewidth}{\wd\pandoc@box}%
  \ifdim\@tempb\p@<\@tempa\p@\let\@tempa\@tempb\fi% select the smaller of both
  \ifdim\@tempa\p@<\p@\scalebox{\@tempa}{\usebox\pandoc@box}%
  \else\usebox{\pandoc@box}%
  \fi%
}
% Set default figure placement to htbp
\def\fps@figure{htbp}
\makeatother
\setlength{\emergencystretch}{3em} % prevent overfull lines
\providecommand{\tightlist}{%
  \setlength{\itemsep}{0pt}\setlength{\parskip}{0pt}}
\usepackage{bookmark}
\IfFileExists{xurl.sty}{\usepackage{xurl}}{} % add URL line breaks if available
\urlstyle{same}
\hypersetup{
  pdftitle={Protocol \textbar{} Duitse Taalvaardigheid},
  pdfauthor={Jesse Postma},
  hidelinks,
  pdfcreator={LaTeX via pandoc}}

\title{Protocol \textbar{} Duitse Taalvaardigheid}
\author{Jesse Postma}
\date{2026-06-05}

\begin{document}
\maketitle

\section{Onderzoeksprotocol}\label{onderzoeksprotocol}

{\def\LTcaptype{none} % do not increment counter
\begin{longtable}[]{@{}
  >{\raggedright\arraybackslash}p{(\linewidth - 2\tabcolsep) * \real{0.2553}}
  >{\raggedright\arraybackslash}p{(\linewidth - 2\tabcolsep) * \real{0.7447}}@{}}
\toprule\noalign{}
\endhead
\bottomrule\noalign{}
\endlastfoot
Onderzoekers & Jesse Postma (507655), Vani Rembet (450024), Dennis
Kuiper (509448) \\
Instelling & Hanze Hogeschool Groningen \\
Datum & Juni 2026 \\
Versie & 1.0 \\
\end{longtable}
}

\subsection{Titel en wetenschappelijke
achtergrond}\label{titel-en-wetenschappelijke-achtergrond}

\begin{center}\rule{0.5\linewidth}{0.5pt}\end{center}

Titel: Duitse taalvaardigheid onder Nederlandse volwassenen.

De beheersing van het Duits staat in Nederland al geruime tijd onder
druk. Tussen 2018 en 2025 daalde het aantal eindexamens Duits met circa
23\% (van \textasciitilde68.000 naar \textasciitilde49.000). Deze
achteruitgrang weerspiegelt zich ook in het hoger onderwijs, waar
universitaire Duits-studies te kampen hebben met dalende
studentelaantallen. Onderzoek suggereert dat dit deels verklaarbaar is
door een curriculum dat zich voornamelijk richt op technische
vaardigheden ten koste van taal- en cultuurspecifieke inhoud. {[}REF?{]}

Tegelijkertijd blijft Duits economisch en geografisch relevant:
Duitsland is de grootste handelspartner van Nederland en de
gemeenschappelijke grens van circa 570 km leidt tot dagelijkse
interactie in grensregio's. Eerder onderzoek {[}@Beerkens2010{]}
{[}@Swarte2013{]} wijst op een mogelijke receptieve meertaligheid in de
grensstreek.

Dit protocol beschrijft de opzet van het surveyonderzoek: van de
meetinstrumenten en wervinsstrategie tot de gegevensverwerking en
overwegingen.

\subsection{Probleemstelling, doelstelling en
onderzoeksvragen}\label{probleemstelling-doelstelling-en-onderzoeksvragen}

\begin{center}\rule{0.5\linewidth}{0.5pt}\end{center}

\subsubsection{Probleemstelling}\label{probleemstelling}

In welke mate beheersen Nederlandse volwassenen de Duitse taal in 2026,
en in hoeverre spelen leeftijd, opleidingsniveau en woonlocatie
(grensregio versus binnenland) daar een rol in?

\subsubsection{Doelstelling}\label{doelstelling}

Inzicht krijgen in de huidige stand van de Duitse taalvaardigheid onder
Nederlandse volwassenen, uitgesplits naar leeftijdsgroep,
opleidingsniveau en nabijheid van de Duitse grens, met behulp van een
gevalideerde taaltoets en achtergrondenquete.

\subsubsection{Hoofdvraag}\label{hoofdvraag}

In hoeverre verschilt de Duitse taalvaardigheid tussen Nederlandse
volwassenen van 40 jaar en ouder en jongere volwassenen (18-39 jaar)?

\subsubsection{Deelvragen}\label{deelvragen}

\begin{itemize}
\tightlist
\item
  Wat is het gemiddelde CEFR-niveau van Duitse taalvaardigheid per
  leeftijdsgroep (18-39 jaar en 40+ jaar)?
\item
  Welke rol speelt de woonlocatie (grensstreek vs.~binnenland) bij de
  actieve beheersing van het Duits?
\item
  Heeft het opleidingsniveau invloed op de beheersing van de Duitse
  taal?
\item
  In welke situaties gebruiken Nederlanders de Duitse taal actief?
\item
  Hoe beoordelen verschillende leeftijdsgroepen het belang van de Duitse
  taal voor Nederland?
\end{itemize}

\subsection{Hypothesen}\label{hypothesen}

\begin{center}\rule{0.5\linewidth}{0.5pt}\end{center}

\subsubsection{Primaire hypothese
(leeftijd)}\label{primaire-hypothese-leeftijd}

\begin{itemize}
\tightlist
\item
  H0: Er is geen significant verschil in het gemiddelde niveau van
  Duitse taalvaardigheid tussen personen van 40 jaar en ouder en
  personen jonger dan 40 jaar.
\item
  H1: Personene van 40 jaar en ouder hebben een hoger gemiddeld niveau
  van Duitse taalvaardigheid dan jongere volwassenen (18-39).
\end{itemize}

\subsubsection{Secundaire hypothesen}\label{secundaire-hypothesen}

\textbf{Locatie} - H0: er is geen veschil in gemiddelde taalbeheersing
tussen personen in de grensstreek en personen in het binnenland. - H1:
Personen die dichter bij de Duitse grens wonen scoren gemiddeld hoger op
de taalvaardigheidtoets.

\textbf{Opleidingsniveau} - H0: Er is geen verschil in gemiddelde
taalbeheersing tussen personen in de grensstreek en personen in het
binnenland. - H1: Een hoger opleidingsniveau gaat samen met een hogere
gemiddelde taalvaardigheid Duits.

Verwachte CEFR-niveauverdeling op basis van literatuur: gemiddeld A2
voor jongere groep; gemiddelde B1 voor 40+ groep. Niveaus C1 en C2
worden als zeldzaam beschouwd in de algemene populatie. {[}ref ERK
niveaus{]}

\subsection{Onderzoeksontwerp}\label{onderzoeksontwerp}

\begin{center}\rule{0.5\linewidth}{0.5pt}\end{center}

\subsubsection{Type onderzoek}\label{type-onderzoek}

Kwantitatief cross-sectioneel surveyonderzoek. Er wordt éénmalig
gemeten; er is geen voor- en nameting. De methode bestaat uit een
gecombineerd instrument van een achtergrondenquête en een
gestandaardiseerde taaltoets (gebaseerd op CEFR-niveaus A1 t/m C1/2).

\subsubsection{Meetinstrument}\label{meetinstrument}

Het instrument is samengesteld in Microsoft Forms en bestaat uit die
opeenvolgende onderelen: - Deel A: Consentvraag (1 vraag) - Deel B:
Achtergrondenquete: persoonskenmerken, blootstelling, gebruik en
attitude (11 vragen) - Deel C: Taalvaardigheidstoets Duits: 50
meerkeuzevragen op CEFR-niveau A1 t/m C2

De volledige invultijd bedraagt naar schatting 10-15 minuten. Deelname
is volledig anoniem: er worden geen namen, sexes, of e-mailaddressen
opgeslagen.

\subsubsection{Taaltoets: opbouw en
CEFR-mapping}\label{taaltoets-opbouw-en-cefr-mapping}

De taaltoets bestaat uit 50 meerkeuzevragen en beslaat grammaticale en
vocabulaire constructies op oplopend niveau. Onderstaande tabel geeft de
globale verdeling per CEFR-niveau weer:

Scores worden omgezet naar een CEFR-niveau door middel van
drempelwaarden die statistisch worden gebaseerd op de uitslag van de
respondenten.

\subsection{Enquetevragen}\label{enquetevragen}

\begin{center}\rule{0.5\linewidth}{0.5pt}\end{center}

\subsection{Deel A: Toestemming}\label{deel-a-toestemming}

{\def\LTcaptype{none} % do not increment counter
\begin{longtable}[]{@{}llll@{}}
\toprule\noalign{}
\endhead
\bottomrule\noalign{}
\endlastfoot
V1 & Wilt u deelnemen aan het onderzoek? & Multiple Choice & Ja/Nee \\
\end{longtable}
}

Toelichting: Bij ``Nee'' wordt de deelnemer bedankt en de enquete
automatisch geeindigd. Alleen respondenten die ``Ja'' antwoorden
vervolgen het formulier.

\subsection{Deel B: Achtergrondenquete}\label{deel-b-achtergrondenquete}

{\def\LTcaptype{none} % do not increment counter
\begin{longtable}[]{@{}
  >{\raggedright\arraybackslash}p{(\linewidth - 6\tabcolsep) * \real{0.0779}}
  >{\raggedright\arraybackslash}p{(\linewidth - 6\tabcolsep) * \real{0.4870}}
  >{\raggedright\arraybackslash}p{(\linewidth - 6\tabcolsep) * \real{0.0779}}
  >{\raggedright\arraybackslash}p{(\linewidth - 6\tabcolsep) * \real{0.3442}}@{}}
\toprule\noalign{}
\begin{minipage}[b]{\linewidth}\raggedright
Nr.
\end{minipage} & \begin{minipage}[b]{\linewidth}\raggedright
Vraag
\end{minipage} & \begin{minipage}[b]{\linewidth}\raggedright
Type
\end{minipage} & \begin{minipage}[b]{\linewidth}\raggedright
Opties
\end{minipage} \\
\midrule\noalign{}
\endhead
\bottomrule\noalign{}
\endlastfoot
V2 & In welke leeftijdscategorie valt u? & MC & 18-39 / 40+ \\
V3 & Wat is het hoogste opleidingsniveau waarvoor u bent geslaagd? & MC
& Basisonderwijs/ VMBO / MBO / HAVO / VWO / HBO / WO \\
V4 & Hoe ver woont u van de Duitse grens? & Open & Numeriek getal \\
V5 & Heeft u ooit in Duitsland gewoond of gewerkt? & MC & Ja/Nee \\
V6 & Heeft u een Duits familielid of kennis? & MC & Ja/Nee \\
V7 & Hoe vaak gebruikt u Duits in het dagelijks leven? & MC &
Nooit/Zelden/SOms/Vaak/Dageljks \\
V8 & In welke situaties gebruikt u Duits? & MC & Werk / Vakantie /
Familie / Media / Booschappen \\
V9 & Heeft u interesse in het leren van een vreemde taal in het
algemeen? & MC & Ja/Nee \\
V10 & Heeft u interesse in het leren van Duits & MC & Ja/Nee \\
V11 & Vindt u dat Duits voor Nederlanders een belangrijke taal is om te
leren? & MC & Ja/Nee \\
V12 & Hoe goed schat u uw eigen Duits in? & Schaal & 1-5 (Slecht -
Vloeiend) \\
\end{longtable}
}

\subsection{Deel C:
Taalvaardigheidtoets}\label{deel-c-taalvaardigheidtoets}

Alle toetsvragen zijn meerkeuze (vier opties, een correct antwoord).
Hieronder een compleet overzicht:

{\def\LTcaptype{none} % do not increment counter
\begin{longtable}[]{@{}
  >{\raggedright\arraybackslash}p{(\linewidth - 4\tabcolsep) * \real{0.0909}}
  >{\raggedright\arraybackslash}p{(\linewidth - 4\tabcolsep) * \real{0.1157}}
  >{\raggedright\arraybackslash}p{(\linewidth - 4\tabcolsep) * \real{0.7851}}@{}}
\toprule\noalign{}
\begin{minipage}[b]{\linewidth}\raggedright
CEFR
\end{minipage} & \begin{minipage}[b]{\linewidth}\raggedright
Question \#
\end{minipage} & \begin{minipage}[b]{\linewidth}\raggedright
Question
\end{minipage} \\
\midrule\noalign{}
\endhead
\bottomrule\noalign{}
\endlastfoot
A1 & 13 & ich \_\_\_ gern pizza. \\
A1 & 14 & \_\_\_ kinder spielen im park. \\
A1 & 15 & die lampe steht \_\_\_ dem tisch. \\
A1 & 16 & ich habe \_\_\_ hund. \\
A1 & 17 & er \_\_\_ zwei schwestern. \\
A1 & 18 & das ist \_\_\_ tasche. \\
A1 & 19 & mein bruder \_\_\_ sehr gut kochen. \\
A1 & 20 & wann \_\_\_ du nach hause? \\
A1 & 21 & ich trinke \_\_\_ kaffee. \\
A1 & 22 & wir \_\_\_ gern musik. \\
A2 & 23 & ich lerne deutsch, \_\_\_ ich in deutschland arbeiten will. \\
A2 & 24 & er \_\_\_ gestern nach berlin gefahren. \\
A2 & 25 & ich gebe \_\_\_ frau die blumen. \\
A2 & 26 & das bild hängt \_\_\_ der wand. \\
A2 & 27 & ich gehe in \_\_\_ schule. \\
A2 & 28 & ich freue \_\_\_ auf den urlaub. \\
A2 & 29 & ich habe \_\_\_ neuen computer gekauft. \\
A2 & 30 & welke zin is juist? \\
A2 & 31 & \_\_\_ sie bitte langsamer! \\
A2 & 32 & das kind hat sich \_\_\_ erkältet. \\
B1 & 33 & welcher satz verwendet den relativsatz richtig? \\
B1 & 34 & wenn ich reich \_\_\_, würde ich um die welt reisen. \\
B1 & 35 & das ist die frau, \_\_\_ ich gestern geholfen habe. \\
B1 & 36 & ich helfe \_\_\_ kollegin gern. \\
B1 & 37 & welke zin is grammaticaal correct? \\
B1 & 38 & es ist wichtig, viel wasser \_\_\_ trinken. \\
B1 & 39 & welke zin is grammaticaal correct?2 \\
B1 & 40 & welke zin drukt een (onrealistische) wens uit? \\
B1 & 41 & ich werde dich morgen \_\_\_. \\
B1 & 42 & welke zin is correct? \\
B2 & 43 & in welke zin is het duitse deelwoord grammaticaal correct
gebruikt? \\
B2 & 44 & welke zin is stylistisch en grammaticaal correct? \\
B2 & 45 & die \_\_\_ der ergebnisse dauert noch an. \\
B2 & 46 & welke zin is grammaticaal correct?3 \\
B2 & 47 & er behauptet, er \_\_\_ nichts davon gewusst. \\
B2 & 48 & welke zin is correct? \\
B2 & 49 & welke zin drukt een vermoeden over het verleden uit? \\
B2 & 50 & in welke zin is een werkwoord correct als zelfstandig
naamwoord gebruikt? \\
B2 & 51 & welke zin is grammaticaal correct?4 \\
B2 & 52 & die frage \_\_\_ noch einer klärung. \\
C1/2 & 53 & die entscheidung \_\_\_ auf heftigen widerstand. \\
C1/2 & 54 & welke zin gebruikt de duitse `gevoelswoordjes' (zoals `halt'
of `eben') op de juiste manier? \\
C1/2 & 55 & welke zin is stylistisch en grammaticaal correct?2 \\
C1/2 & 56 & welke zin drukt een tegenstelling (met `hoewel')
grammaticaal correct uit? \\
C1/2 & 57 & in welke zin wordt de uitspraak van de minister op de
grammaticaal juiste manier naverteld? \\
C1/2 & 58 & welke zin is grammaticaal correct?5 \\
C1/2 & 59 & die studie lässt \_\_\_ schließen, dass weitere forschung
nötig ist. \\
C1/2 & 60 & welke zin bevat de juiste bijzin? \\
C1/2 & 61 & der bericht, \_\_\_ ergebnisse kontrovers diskutiert wurden,
liegt nun vor. \\
C1/2 & 62 & es \_\_\_ noch zu klären, inwieweit diese regelung anwendbar
ist. \\
\end{longtable}
}

\subsection{Doelgroep, inclusie- en
exclusiecriteria}\label{doelgroep-inclusie--en-exclusiecriteria}

\begin{center}\rule{0.5\linewidth}{0.5pt}\end{center}

\subsubsection{Doelgroep}\label{doelgroep}

Nederlandse volwassenen van 18 jaar en ouder, zonder Duits als
moedertaal, die tenminste tien jaar in Nederland wonen en het Nederlands
voldoende beheersen om een enquete te begrijpen.

\subsubsection{Inclusiecriteria}\label{inclusiecriteria}

\begin{itemize}
\tightlist
\item
  Leeftijd 18 jaar of ouder
\item
  Woont minimaal 10 jaar in Nederland
\item
  Duits is niet een moedertaal
\item
  Voldoende Nederlands om de vragenlijst in te vullen
\item
  Bereid tot deelname na mondelinge toestemming
\end{itemize}

\subsubsection{Exclusiecriteria}\label{exclusiecriteria}

\begin{itemize}
\tightlist
\item
  Leeftijd jonger dan 18 jaar
\item
  Duits als moedertaal (e.g.~Duitse studenten/expats in Nederland)
\item
  Minder dan 10 jaar wonend in Nederland
\item
  Onvoldoende begrip Nederlandse taal
\item
  Geen consent voor de vragenlijst
\end{itemize}

\subsubsection{Steekproefomvang en
-methode}\label{steekproefomvang-en--methode}

Nagestreefde minimale steekproef: 40 respondenten. Uiteindelijk
gerealiseerd: 154 respondenten, weergegeven in de dataset(s) in
`/ruwe\_data'.

Steekproefmethode: gelegenheidssteekproef, aangevuld door werving via
persoonlijk netwerk. Gekozen gezien gemakkelijke wegen van connectie via
bijvoorbeeld social media. Daarnaast waren er flyers beschikbaar voor
gemakkelijke toegang tot de vragenlijst in een openbare ruimte.

\subsection{Wervingsstrategie}\label{wervingsstrategie}

\begin{center}\rule{0.5\linewidth}{0.5pt}\end{center}

\subsection{Kanalen}\label{kanalen}

Respondenten werden via de volgende kanalen geworven:

{\def\LTcaptype{none} % do not increment counter
\begin{longtable}[]{@{}
  >{\raggedright\arraybackslash}p{(\linewidth - 4\tabcolsep) * \real{0.1438}}
  >{\raggedright\arraybackslash}p{(\linewidth - 4\tabcolsep) * \real{0.4118}}
  >{\raggedright\arraybackslash}p{(\linewidth - 4\tabcolsep) * \real{0.4379}}@{}}
\toprule\noalign{}
\begin{minipage}[b]{\linewidth}\raggedright
Kanaal
\end{minipage} & \begin{minipage}[b]{\linewidth}\raggedright
Methode
\end{minipage} & \begin{minipage}[b]{\linewidth}\raggedright
Opmerkingen
\end{minipage} \\
\midrule\noalign{}
\endhead
\bottomrule\noalign{}
\endlastfoot
Openbare ruimtes & Benadering met QR-code (poster). & Werk, cafe's \\
Sociaal netwerk & WhatsApp, mondelinge verzoeken aan
vrienden/familie/kenissen & Meeste respondenten afkomstig uit sociaal
netwerk \\
Facebook, Instagram & Posts met aanroepen tot invullen & Gerichte posts
in enquetegroepen/ 40+ groepen ook van toepassing \\
SurveyCircle & ``Last measure'' & 0 respondenten geworven \\
\end{longtable}
}

\subsubsection{Mondelinge
wervingsscript}\label{mondelinge-wervingsscript}

Bij persoonlijke benadering in openbare ruimtes werd het volgende script
gehanteerd:

\textbf{Stap 1 - opening:}

``Hallo, wij doen voor onze studie een onderzoek naar hoe goed
Nederlanders nog Duits kunnen. Zou u ons willen helpen? Het duurt
ongeveer x minuten.''

\textbf{Stap 2 - Screening taal:}

``Is Duits uw moedertaal?'' - Ja -\textgreater{} ``Dan valt u buiten
onze doelgroep. Bedankt voor uw tijd.'' - Nee -\textgreater{}
voortzetting stap 3

\textbf{Stap 3 - Screening woonsituatie:}

``Hoe lang woont u al in Nederland?'' - antwoord \textgreater{} 10
-\textgreater{} voortzetting stap 4 - antwoord \textless{} 10 jaar
-\textgreater{} ``Dan valt u buiten onze doelgroep. Bedankt voor uw
tijd.''

\textbf{Stap 4 - Screening leeftijd:}

``Mag ik vragen hoe oud u bent?'' - antwoord \textgreater{} 18
-\textgreater{} voortzetting stap 5 - antwoord \textless{} 18 jaar
-\textgreater{} ``Dan valt u buiten onze doelgroep. Bedankt voor uw
tijd.''

\textbf{Stap 5 - Doelinformatie en overhandiging survey:}

``Wij onderzoeken het verschil in taalbeheersing van verschillende
generaties. Dit doen wij voor een onderzoek bij de Hanze Hogeschool
Groningen. De enquête is volledig anoniem. Wilt u meedoen?‟ - Ja
-\textgreater{} QR-code tonen of toegang tot survey verlenen via ander
medium. - Nee -\textgreater{}''Geen probleem, fijne dag!''

\subsection{Uitvoeringsprocedure}\label{uitvoeringsprocedure}

\begin{center}\rule{0.5\linewidth}{0.5pt}\end{center}

\subsubsection{Materialen}\label{materialen}

\begin{itemize}
\tightlist
\item
  Microsoft Forms-formulier (online)
\item
  A5-poster met duidelijke QR-code en korte insructie.
\item
  Eigen smartapparaat als back-up voor respondenten zonder smartphone.
\item
  Programmeertaal R voor statistische verwerking en visualisatie
\end{itemize}

\subsubsection{Gegevensverzameling}\label{gegevensverzameling}

Elke respondent doorloopt het formulier in drie onderdelen:

\begin{enumerate}
\def\labelenumi{\arabic{enumi}.}
\tightlist
\item
  Toestemmingsvraag
\item
  Achtergrondenquete
\item
  Taaltoets
\end{enumerate}

Alle antwoorden worden automatisch opgeslagen in Microsoft Forms en zijn
exporteerbaar als .xlsx. De dataset is geconverteerd naar een .csv
bestand voor analyse in R.

\subsubsection{Tijdspad}\label{tijdspad}

{\def\LTcaptype{none} % do not increment counter
\begin{longtable}[]{@{}
  >{\raggedright\arraybackslash}p{(\linewidth - 4\tabcolsep) * \real{0.1584}}
  >{\raggedright\arraybackslash}p{(\linewidth - 4\tabcolsep) * \real{0.2376}}
  >{\raggedright\arraybackslash}p{(\linewidth - 4\tabcolsep) * \real{0.5941}}@{}}
\toprule\noalign{}
\begin{minipage}[b]{\linewidth}\raggedright
Week
\end{minipage} & \begin{minipage}[b]{\linewidth}\raggedright
Periode
\end{minipage} & \begin{minipage}[b]{\linewidth}\raggedright
Activiteit
\end{minipage} \\
\midrule\noalign{}
\endhead
\bottomrule\noalign{}
\endlastfoot
1 & 20-26 april 2026 & Onderwerpsselectie en literatuurverkenning. \\
2 & 27 april - 3 mei 2026 & Docentsfeedback; opstelling PvA \\
3 & 4-10 mei 2026 & Bouwen Forms, drukken posters, start werving \\
4 & 11-17 mei & Dataverzameling \\
5 & 18-24 mei & Data-export, opschonen data, beginnen exploratie binnen
R \\
6 & 25-31 mei & Statistische analyses en visualisaties in R \\
7 & 1-7 juni & Schrijven paper, verdere analyse en inhoud \\
8 & 7-9 juni & Eindcorrectie en inlevering \\
\end{longtable}
}

\subsection{Ethische overwegingen}\label{ethische-overwegingen}

\begin{center}\rule{0.5\linewidth}{0.5pt}\end{center}

\subsubsection{Anonimiteit en informed
consent}\label{anonimiteit-en-informed-consent}

Deelname aan het onderzoek is volledig anoniem. Er worden geen
persoonsidentificerende gegevens opgeslagen (geen naam, geen
e-mailadres, geen IP-adres). De eerste vraag van het formulier is een
expliciete toestemmingsvraag; respondenten die ``Nee'' antwoorden worden
direct doorgestuurd naar een bedankpagina en hun (lege) respons wordt
niet opgeslagen.

\subsubsection{Kwetsbare groepen}\label{kwetsbare-groepen}

Er is geen doelgerichte werving onder kwetsbare groepen. Minderjarigen
(jonger dan 18 jaar) worden als inclusievraag uitgesloten. Bij
benadering in openbare ruimtes werden potentiële respondenten die
tekenen van cognitieve beperkingen vertoonden niet actief benaderd.

\subsubsection{Dataveiligheid}\label{dataveiligheid}

De ruwe data worden opgeslagen op een GitHub-repository. Toegang is
openbaar, al is de data volledig geanonimiseerd. Bij publicatie of
inlevering wordt de dataset korte tijd bewaard conform de richtlijnen
voor studentenonderzoek aan de Hanze Hogeschool Groningen.

\subsection{Risicoanalyse}\label{risicoanalyse}

\begin{center}\rule{0.5\linewidth}{0.5pt}\end{center}

{\def\LTcaptype{none} % do not increment counter
\begin{longtable}[]{@{}
  >{\raggedright\arraybackslash}p{(\linewidth - 4\tabcolsep) * \real{0.3350}}
  >{\raggedright\arraybackslash}p{(\linewidth - 4\tabcolsep) * \real{0.2750}}
  >{\raggedright\arraybackslash}p{(\linewidth - 4\tabcolsep) * \real{0.3850}}@{}}
\toprule\noalign{}
\begin{minipage}[b]{\linewidth}\raggedright
Risico
\end{minipage} & \begin{minipage}[b]{\linewidth}\raggedright
Impact
\end{minipage} & \begin{minipage}[b]{\linewidth}\raggedright
Mitigatie
\end{minipage} \\
\midrule\noalign{}
\endhead
\bottomrule\noalign{}
\endlastfoot
Respondenten buiten doelgroep (toeristen, expats, minderjarigen) &
Vertekening resultaten & Screening op basis van leeftijd, moedertaal,
verblijfsduur. \\
Selectiebias:

oververtegenwoordiging sociale kring & Beperkte externe validiteit,
geografische clustering & Actieve werving in openbare ruimtes en via
netwerken als Facebook. \\
Lage respons oudere groepen & Oneven verdeling leeftijdsgropen &
Gerichte benadering \\
Begripsproblemen toetsvragen & Meting van begrijpend lezen ipv
taalvaardigheid & Korte instructieteksten, simpele uitleg, meerkeuze
format. \\
Onnauwkeurige inschatting & Wenselijkheidbias & Vergelijken schatting vs
werkelijke score \\
Surveycircle (0 respondents) & Geen impact; kanaal niet representatief &
- \\
Onvolledige invulling & Gemiste/onvolledige data & Pauzeren forms
mogelijk, lokken met zien uiteindelijke score op het einde. \\
\end{longtable}
}

\subsection{Gegevensverwerking en statistische
analyse}\label{gegevensverwerking-en-statistische-analyse}

\begin{center}\rule{0.5\linewidth}{0.5pt}\end{center}

\subsubsection{Datastructuur}\label{datastructuur}

De dataset wordt geexporteerd vanuit Microsoft Forms als .xlsx en
omgezet naar .csv. Elke rij representeert een respondent (n = 154).
Kolommen zijn onderverdeeld in: - Administratieve kolommen (ID,
tijdstempels, Forms-metadata)\\
- Responskolommen achtergrond - Responskolommen taaltoets

\subsubsection{Geplande analyses in R}\label{geplande-analyses-in-r}

\begin{itemize}
\tightlist
\item
  Beschrijven resultaten
\item
  Hypothese testen
\item
  Exploratie andere data
\end{itemize}

\subsubsection{Statistische Methoden}\label{statistische-methoden}

• Beschrijvende statistiek: gemiddelden, medianen, standaarddeviaties
per groep • Groepsvergelijking: onafhankelijke t-toets (twee
leeftijdsgroepen) of Mann-Whitney U (bij niet-normale verdeling) •
Correlatie: Spearmans rangcorrelatie tussen afstand tot de grens en
toetsscore • Variantieanalyse: ANOVA of Kruskal-Wallis voor vergelijking
drie of meer opleidingsniveaus

\subsection{Fair-principes en
management}\label{fair-principes-en-management}

\subsubsection{Findable}\label{findable}

• Dataset heeft een beschrijvende bestandsnaam: resultaat.xlsx /
resultaat.csv • Een README-bestand beschrijft de inhoud van elke kolom •
Mappenstructuur: /ruwe\_data, /analyse, /figures, /scripts

\subsubsection{Accessible}\label{accessible}

• Opgeslagen op Hanze OneDrive (toegankelijk voor onderzoekers en
begeleiders) • Geback-upt in privé GitHub-repository • Bestanden te
openen in Excel, R en LibreOffice Calc zonder speciale software

\subsubsection{Interoperable}\label{interoperable}

• Export als open formaat (.csv) voor gebruik in R • Kolomnamen zijn
beschrijvend en gestandaardiseerd • Aandachtspunt: één cel bevat
meerdere waarden (V8: meervoudig mc) --- vereist opschoning in R
(one-hot encoding).

\subsubsection{Reusable}\label{reusable}

• Dataverzamelingsprotocol (dit document) beschrijft alle stappen
volledig • Inclusie- en exclusiecriteria zijn gedocumenteerd •
Hergebruik door een externe onderzoeker is mogelijk mits toegang tot de
README

\subsection{Literatuur}\label{literatuur}

Abitzsch, D., Knopp, E., \& Michel, M. (2024). Motivatie om Duits te
studeren aan Nederlandse universiteiten. Levende Talen Tijdschrift, (2),
13--24. Beeker, A., Fasoglio, D., de Jong, K., Keuning, J., \& van Til,
A. (2015). Weergave van ERK-niveau schrijfvaardigheid Engels, Duits en
Frans. Levende Talen Tijdschrift, (2), 16--25. Beerkens, R. (2010).
Receptive multilingualism as a language mode in the Dutch-German border
area. Waxmann Publishing Co. Gille, E., \& Kordes, J. (2013). Weergave
van vaardigheid Engels en Duits van Nederlandse leerlingen in Europees
perspectief. Levende Talen Magazine, (1), 23--27. Swarte, F., Schüppert,
A., \& Gooskens, C. (2013). Do speakers of Dutch use their knowledge of
German while processing written Danish words? Linguistics in the
Netherlands, 30(1), 146--159. Van den Noort, M. W. M. L., Bosch, P., \&
Hugdahl, K. (2006). Foreign language proficiency and working memory
capacity. European Psychologist, 11(4), 289--296. Verkuil, D. (2015,
maart). Hoge prijs voor gebrekkig Duits. NOS.
\url{https://nos.nl/artikel/2026034} Weltens, B., de Bot, K., \& van
Els, T. (1986). Language attrition in progress. Floris Publications.
Wilhelm, F. (2018). Foreign language teaching and learning in the
Netherlands 1500--2000. The Language Learning Journal, 46(1), 17--27.
Zinecker, M., \& Konečná, Z. (2014). German language proficiency among
students of business and management in the Czech Republic. DANUBE: Law
and Economics Review, 5(3), 173--200.

\end{document}
